%%%%%%%%%%%%%%%%%%%%%%%%%%%%%%%%%%%%%%%%%%%%%%%%%%%%%%%%%%%%%%%%%%
\section{Introduction}
\label{sec:intro}
%%%%%%%%%%%%%%%%%%%%%%%%%%%%%%%%%%%%%%%%%%%%%%%%%%%%%%%%%%%%%%%%%%

\noindent
Mooring lines, risers, aquaculture netting, textile yarns, and aquatic vegetation all have a cross-section orders of magnitude smaller than the flow domain they occupy, and they all deform under the flow that loads them. Resolving such a body with a boundary-fitted mesh imposes severe local refinement, and the mesh must then be deformed or regenerated at every time step as the structure moves. The computational cost becomes prohibitive for a canopy of vegetation stems or a net of many strands.

\noindent
Distributed forcing methods avoid that cost by representing the structure through a lower-dimensional description and introducing its action on the fluid as regularised momentum sources on a fixed Eulerian mesh. The immersed boundary method \citep{peskin_flow_1972,peskin_immersed_2002,mittal_immersed_2005,uhlmann_immersed_2005} established the idea for surfaces, and the actuator line method of \citet{sorensen_numerical_2002} reduced it further for slender members by concentrating the hydrodynamic interaction at actuator points along a curve. Developed for wind-turbine wakes, where it captures the momentum exchange on meshes far coarser than blade-resolved calculations \citep{troldborg_numerical_2010}, it has since been refined through improved sampling, smoothing, and generalised actuator-curve representations \citep{churchfield_advanced_2017,jha_actuator_2018,navarro_diaz_actuator_2023}. Nothing in the concept restricts it to rotating blades. \citet{bral_development_2025} adapted it to flexible yarns in air-jet weaving and then extended it into a two-way coupled framework driving a beam model \citep{bral_modeling_2025}, while \citet{schmitt_coupled_2022} coupled an actuator line to a finite element solver in OpenFOAM for flexible turbine blades and fish-farm nets.

\noindent
Two operators govern the accuracy of such a coupling, namely velocity sampling and force regularisation. Sampling may be done at a point, along a line, or as a volume average \citep{sorensen_numerical_2002,merabet_parametric_2019,zormpa_effect_2025}. The choice matters most for drag-dominated members, because the regularised source acts as a local sink that slows the fluid in the same region where the velocity would be sampled, and the drag scales with the square of that velocity, so an underestimate feeds back on itself. \citet{bral_development_2025} broke this feedback by sampling at a fixed distance upstream of the actuator point, outside the momentum-deficit region, and the same concern underlies the sampling and smearing corrections of \citet{churchfield_modeling_2015}, \citet{martinez_tossas_optimal_2017} and \citet{muscari_effective_2024}. Regularisation is needed because a singular point force cannot be applied to any discretisation of finite resolution, whether a finite volume or finite element mesh, a finite difference grid, or a set of particles, without a severe dependence on that resolution. The Gaussian kernel of width $\varepsilon$ that replaces it leaves the source grid-localised when too narrow and weakens the local effect of the line when too wide \citep{martinez_tossas_optimal_2017,navarro_diaz_actuator_2023}. The two operators also fix which discrete properties the coupling preserves, namely the conservation of exchanged momentum, the resolution of the kernel relative to the local grid, and invariance under parallel domain decomposition \citep{bral_development_2025}.

\noindent
On the structural side, the geometrically exact beam theory of \citet{reissner_finite_1981}, extended to three dimensions by Simo, treats finite rotations without linearisation and retains shear and torsion alongside bending and axial stretch. Most implementations of it are finite element, and in coupled simulations the structure is usually solved by a different numerical method from the flow. Numerical studies of flexible vegetation illustrate this, with the stems represented by finite element elastic rods \citep{chen_eulerianlagrangian_2019,elrahi_numerical_2023}, immersed-boundary soft bodies \citep{kim_steady_2024}, direct-forcing immersed-boundary solvers \citep{pruter_comprehensive_2025}, or articulated-body methods \citep{dickinson_modelling_2025}, each coupled to a separate flow solver. The two sides then hold their data differently, are decomposed differently across processors, and must each be verified in their own setting first. One way to remove that separation is to discretise the flow and the structure within the same numerical framework, whether finite element or finite volume. The present work takes the finite volume route, using the cell-centred finite volume discretisation of the same beam theory developed by \citet{bali_cell-centered_2022,bali_finite_2024,bali_beamfoam_2025}, so that the two subsystems share one set of conservation statements, one mesh and field infrastructure, one domain decomposition, and one time loop. Two-way actuator-line coupling has been demonstrated for flexible marine structures, but on finite element or external structural solvers \citep{schmitt_coupled_2022,bral_modeling_2025}, so the sampling, regularisation, and parallel exchange have not been established for a geometrically exact finite volume beam.

\noindent
The hydrodynamic closure adopted here is a quasi-steady Morison-type drag law with no added-mass or history term, and a closure of this kind has an applicability boundary that can be stated in advance. \citet{leclercq_reconfiguration_2018} showed that the wave-induced response of a flexible blade is governed by the combination $CaL/KC$, and that below a threshold of approximately $0.2$ the blade responds quasi-statically in a single spatial mode, while above it the added mass of the surrounding water amplifies the motion and excites higher modes. \citet{jacobsen_waveinduced_2019} confirmed this experimentally for single vegetation mimics and derived drag coefficients by fitting a relative-velocity Morison formulation to the measured base force. Those fitted coefficients carry a documented spread between repeat tests, so a validation that fixes the coefficient in advance should propagate that band through the simulation rather than tune a single value against the response being validated.

\noindent
This paper presents a finite volume actuator-line formulation for the fluid--structure interaction of flexible slender structures, in which a cell-centred finite volume Simo--Reissner beam is coupled two-way to an incompressible finite volume flow solver through conservative and parallel-consistent transfer operators, and in which the accuracy limits of a drag-only Morison closure are quantified against benchmark and wave-flume data. No external structural code is required, because the beam shares the mesh and field infrastructure, the time loop, and the domain decomposition of the flow solver. The sampling and projection operators are formulated to be parallel consistent on a domain-decomposed fluid mesh, and the projection is normalised so that it conserves the transferred resultant force independently of the mesh, the kernel width, and the kernel truncation. The formulation does not depend on the number of fluid phases, and is used unchanged in single-phase channel flow and in a two-phase wave flume. It is verified and validated against three reference problems of increasing difficulty, with every drag coefficient fixed in advance, the two operator settings with a physical rather than a purely numerical influence on the response identified, and the accuracy of the closure together with the range of conditions over which it holds quantified rather than assumed. Verification and validation are used throughout in the standard sense \citep{roache_quantification_1997,oberkampf_roy_2010}, the first established against independent numerical solutions and the second against experiment.

\noindent
Section~\ref{sec:mathmodel} presents the governing equations and the hydrodynamic closure, and Section~\ref{sec:numericalmodel} the finite volume beam discretisation and the coupling operators. Sections~\ref{sec:cantilever} to~\ref{sec:leaflet} verify the structural solver and then the coupling, against a flexible beam in a steady confined channel and a leaflet in a sinusoidal channel flow. Section~\ref{sec:flexveg} validates the coupling against the wave-flume experiment of \citet{jacobsen_waveinduced_2019} on a flexible vegetation stem, and Section~\ref{sec:conclusion} draws the conclusions.
